% =============================================================================
% Efficiency / Cost Comparison Table — for ESWA Rebuttal / Camera-Ready
% Dataset: alpaca_eval (n = 161)
% Metrics:
%   - AQSA  : Average Queries per Successful Attack
%   - Tokens/Success (M) : per-sample token consumption until first success, in millions
%   - GPU time: end-to-end wall-clock time on a single GPU
% Baselines: GCG, AutoDAN, TAP
% =============================================================================
% Placement suggestion: insert in the Experiments section, immediately after
% the main ASR comparison table (so readers see "effectiveness" then "cost").
% Or in a dedicated \subsection{Computational Efficiency}.
% -----------------------------------------------------------------------------

\paragraph{Computational efficiency.}
Beyond attack effectiveness, we compare PreferAttack against three
strong baselines — GCG~\cite{zou2023gcg}, AutoDAN~\cite{liu2024autodan},
and TAP~\cite{mehrotra2023tap} — on three efficiency dimensions:
(i) \emph{average queries per successful attack} (AQSA),
(ii) \emph{per-sample token consumption until first success} (in
millions, M), and (iii) \emph{total GPU wall-clock time} on the full
alpaca\_eval test split ($n{=}161$).
All methods are evaluated under the same judge configuration
(Qwen-2.5-1.5B-Instruct) on a single NVIDIA [GPU-TYPE-PLACEHOLDER]
GPU with batch size $64$. Results are summarised in
Table~\ref{tab:efficiency}.

% -----------------------------------------------------------------------------
% Suggested label: tab:efficiency
% Suggested caption: Efficiency comparison on alpaca_eval (n = 161).
% -----------------------------------------------------------------------------
\begin{table}[t]
  \centering
  \caption{Efficiency comparison on \textbf{alpaca\_eval} ($n{=}161$).
  \textbf{AQSA}: average queries per successful attack (\emph{lower is better}).
  \textbf{Tokens / Success}: per-sample token consumption until first success,
  in millions (\emph{lower is better}).
  \textbf{GPU Time}: end-to-end wall-clock time on a single NVIDIA
  [GPU-TYPE-PLACEHOLDER] GPU (\emph{lower is better}).
  Bold indicates the best trade-off across the three metrics.}
  \label{tab:efficiency}
  \small
  \begin{tabular}{lccc}
    \toprule
    \textbf{Method}
      & \textbf{AQSA} $\downarrow$
      & \textbf{Tokens / Success (M)} $\downarrow$
      & \textbf{GPU Time} $\downarrow$ \\
    \midrule
    GCG              & $60{,}184.42$          & $24.80$           & $2.71$\,h          \\
    AutoDAN          & $204.90$               & $0.49$            & $1$\,h\,$43$\,min  \\
    TAP              & $\mathbf{9.93}$        & $\mathbf{0.027}$  & $3.10$\,h          \\
    \textbf{Ours}    & $\mathbf{34.40}$       & $\mathbf{0.80}$   & $\mathbf{1.50}$\,\textbf{h} \\
    \bottomrule
  \end{tabular}
\end{table}

% -----------------------------------------------------------------------------
% Discussion paragraph — interpret the table.
% -----------------------------------------------------------------------------
The three baselines exhibit qualitatively different cost profiles:
\begin{itemize}[leftmargin=*]
  \item \textbf{GCG} is by far the most token-hungry baseline: its
        token-level gradient optimisation requires $\sim\!60\text{K}$
        queries and $\sim\!24.8$\,M tokens per successful attack --- two
        to three orders of magnitude more than any other method --- while
        still finishing in $2.71$\,h due to its batched implementation.
  \item \textbf{AutoDAN} replaces gradient-based optimisation with a
        genetic algorithm, cutting both queries ($\sim\!205$) and
        per-sample tokens ($0.49$\,M) by $\sim\!2$--$3$ orders of
        magnitude relative to GCG. It is the fastest in wall-clock time
        ($1$\,h\,$43$\,min), but its ASR is materially lower than the
        LLM-driven methods (cf.\ Table~\ref{tab:multiseed}).
  \item \textbf{TAP} achieves the best AQSA ($9.93$) and the lowest
        per-sample token cost ($0.027$\,M) thanks to its LLM-based
        attacker that rewrites instructions in a single pass. However,
        each LLM generation step is expensive, making it the
        \emph{slowest} method overall at $3.10$\,h --- roughly
        $2\times$ our wall-clock time.
  \item \textbf{PreferAttack (Ours)} strikes a favourable balance:
        $34.4$ queries per success, $0.80$\,M tokens, and
        $1.5$\,h GPU time --- only $\sim\!5\%$ slower than AutoDAN
        while delivering substantially higher ASR, and $\sim\!2\times$
        faster than TAP with comparable query efficiency.
\end{itemize}

In short, PreferAttack achieves the best \emph{effectiveness--efficiency
trade-off}: its per-sample attack cost is two orders of magnitude lower
than GCG and within $2\times$ of TAP (the most token-efficient
baseline), while its wall-clock time is competitive with the fastest
baseline (AutoDAN) and nearly half of TAP's. The multi-agent
decomposition --- with the GA doing heavy exploration in token space
and the RL controller dynamically reallocating search effort ---
allows PreferAttack to converge to successful suffixes without
incurring the per-step LLM-generation cost of TAP or the brute-force
token-level search cost of GCG.

% -----------------------------------------------------------------------------
% NOTE FOR AUTHORS — fill these placeholders before submission:
%   1. [GPU-TYPE-PLACEHOLDER]  — replace with the actual GPU used (e.g., A100-80G, H100, RTX 4090)
%   2. \cite{...} keys         — verify they match your .bib file
%   3. 1h43min for AutoDAN     — your input was "1h26min"; please double-check.
%       The number 1.43\,h corresponds to 1h26min. If the actual measured
%       time differs, update both the table cell and the discussion text.
%   4. Cross-reference         — Table~\ref{tab:multiseed} should point to
%                                your main ASR table; adjust label if needed.
% -----------------------------------------------------------------------------
